\section{Discussion}\label{sec:discussion}

\para{Hypotheses revisited.}
Our five hypotheses are borne out to differing degrees, and reading them together tells a
coherent story. {\bf Does sampling cost accuracy (H1)?} Yes: single-sample accuracy is
$\leq$ \detmode{} for every genuine sampling variant, and the gap is not a fixed tax but
grows monotonically as we sample harder. Sharpening the categorical sample (lowering $\tau$)
drives MNIST accuracy from $97.6\%$ down to $66.4\%$ (\Tabref{tab:e2}), and moving the
distribution from continuous to binary to categorical widens the cost from $\approx 0$ to
$\approx 2.8$ to $\approx 10$ points (\Tabref{tab:e1}). {\bf Does sampling regularize
(H2)?} Yes: every stochastic variant shrinks the train--test gap, and most lower \ece{}.
\gaussreparam{} and \bernoullist{} improve calibration \emph{and}, on \cifar{}, \nll{}
(\gaussreparam{} cuts NLL $0.94\to0.78$ and ECE $0.112\to0.074$; \tabref{tab:e6}). {\bf Does
MC-averaging recover the cost (H3)?} Yes (\figref{fig:e4}): averaging many stochastic passes
approximates the expectation over the activation distribution, exactly as averaging softmax
draws approximates the class distribution, lifting \dropout{} from $96.6\%$ to $97.9\%$ as
$K$ grows $1\to50$. {\bf Does the distribution type govern the outcome (H4)?} Strongly: the
Gaussian--Bernoulli--categorical ordering is visible in accuracy, in \nll{}, and in gradient
variance, and the gradient-variance ranking (\gaussnoise{}/\bernoullist{} $\approx10^{-7}$
$<$ \dropout{}/\gaussreparam{} $\approx10^{-6}$ $<$ \gumbelsoft{} $7\times10^{-6}$ $<$
\gumbelst{} $1.5\times10^{-5}$) mechanistically \emph{explains} the accuracy ranking. {\bf
Does sampling improve robustness (H5)?} Only weakly: the benefit under input noise is modest
and comes almost entirely from \dropout{} with MC-averaging ($+1$ point at $\sigma=1.0$,
$+2.2$ at $\sigma=1.5$; \figref{fig:e5}), not from the sampling variants themselves.

\para{Why categorical sampling is special.}
The four sampling operators are not equally severe, and the reason is sharp: categorical
sampling is the only one that changes the \emph{dimensionality} of information flow. A hard
one-hot draw from a softmax over $W$ hidden units passes essentially $\log_2(W)$ bits
forward, whereas \gaussreparam{} and \bernoullist{} pass a full per-unit vector. This single
fact is simultaneously why categorical sampling is the \emph{strongest} regularizer --- the
smallest generalization gap on MNIST ($-0.021$ for \gumbelst{}) and the lowest \ece{} on
\cifar{} --- and why it is the highest-variance, hardest-to-train operator. The same
bottleneck that prevents the network from memorizing also prevents it from carrying enough
information when the task needs it. On MNIST, where the decision is nearly linearly
separable, a mild relaxation (large $\tau$) is a tolerable squeeze; on \cifar{}, forcing all
CNN features through one sampled unit is too tight, and training collapses to chance
(\tabref{tab:e6}). This is precisely the regime the information bottleneck literature
describes \citep{tishby2015deep,alemi2017deepvib}: compressing the representation trades
predictive sufficiency for minimality, and pushed too far the representation becomes
insufficient. It is also a concrete instance of the codebook-collapse caution familiar from
discrete latent models \citep{oord2017vqvae}, where a categorical bottleneck degenerates to a
handful of active codes. The \detmode{}, \gaussreparam{}, and \bernoullist{} operators never
touch this dimensionality and never collapse.

\para{Surprises.}
Three results ran against our priors. First, both Gaussian variants --- additive
\gaussnoise{} and reparameterized \gaussreparam{} --- are effectively \emph{free}
($97.7\%$ vs $97.8\%$ on MNIST, $p=0.83$ for \gaussreparam{}), and on \cifar{}
\gaussreparam{} improves accuracy, calibration, \emph{and} \nll{} together ($75.5\%$ vs
$75.0\%$, ECE $0.074$ vs $0.112$, NLL $0.78$ vs $0.94$). A continuous sample placed at a
hidden layer is not a cost to be paid but, in this setting, a small net win. Second, the
\gumbelsoft{} (soft) variant collapses to \emph{exactly} chance ($10.0\%$) on \cifar{} ---
not merely degraded but untrainable, which we did not anticipate from a relaxed, fully
differentiable operator. Third, sampling \emph{later} is cheaper (\tabref{tab:e3}): moving
the \samplinglayer{} nearer the output lifts \bernoullist{} from $95.0\%$ to $97.8\%$, the
opposite of a naive ``squeeze early to force a good bottleneck'' intuition. The bottleneck
hurts least when it sits nearest the already-low-dimensional decision.

\para{Limitations.}
Several caveats bound these conclusions. We study two datasets (\mnist{}, \cifar{}), small
models, a single sampling site per network, and 3 seeds; the trends are clear and
significant, but absolute numbers are modest, with no heavy CIFAR tuning or augmentation. The
most important caveat concerns the categorical collapse: it could be partly an
\emph{optimization} failure --- the learning rate, scale, or temperature for the
width-rescaled one-hot at a wide fully-connected layer --- rather than a pure information
limit. A $\tau$-annealing schedule and per-mode learning-rate tuning were out of scope; we
report the collapse honestly under the shared protocol rather than engineering it away, and
the monotonic $\tau$-sweep (\figref{fig:e2}) supports the information-bottleneck reading that
the effect is real and graded. We also note that MC-averaging \emph{categorical} samples can
\emph{raise} \ece{}, because averaging near-one-hot draws yields overconfident mean
probabilities --- MC-averaging is not a universal calibration fix. Robustness was tested only
under Gaussian input noise, not adversarial or corruption suites. Finally, gradient variance
was measured on one parameter tensor feeding the sampler as a proxy for the operator's
overall noise, not as an exhaustive accounting.

\para{Broader implications.}
The practical guidance is asymmetric and, we think, useful. Continuous (Gaussian)
intermediate-layer sampling is a near-free drop-in: \gaussreparam{} costs no measurable
accuracy and reliably improves calibration and the generalization gap, making it a sensible
default when better-behaved probabilities are wanted at little risk. The literal categorical
idea --- ``sample one hidden unit the way we sample the output softmax'' --- is \emph{not} a
free drop-in. It is a severe information bottleneck whose cost scales with three things at
once: task difficulty (gentle on MNIST, catastrophic on \cifar{}), sampling sharpness
(monotonic in $\tau$), and earlier placement (cheaper near the output). Practitioners drawn
to the elegant softmax-sampling analogy should treat categorical hidden sampling as a
deliberate, heavily-regularizing bottleneck to be annealed and tuned, not as a plug-in
stochastic activation. More broadly, our controlled comparison suggests that the right
question about hidden-layer stochasticity is not \emph{whether} to sample but \emph{from what
distribution}: the choice of distribution type, far more than the mere act of sampling,
determines whether the operation is a free regularizer or a training-breaking bottleneck.
